\section{铁路转弯处的过渡曲线}
\markboth{86. 铁路转弯处的过渡曲线}{86. 铁路转弯处的过渡曲线}

在原野上奔驰着的火车，转弯的时候，为什么不会翻车呢？学过中学物理的人都知道，这是因为在铺设铁轨之前已经作过精密的计算，根据火车拐弯时所需向心力的大小，事先已把转弯处的外侧轨道填高了，所以火车转弯时万无一失。

然而这里的问题还很不简单。火车在直道上行驶时，两根平行的轨道是一样的高低。但是改在弯道上行驶时，外轨却要比内轨高，其高度差 $h$ 与轨道半径 $R$ 有关。从两根轨道一样高，改变到一根比另一根高 $h$，必须经历一个过渡阶段，让两根轨道的高度差从 $0$ 逐渐连续地增大到 $h$。在这过渡阶段中，由于车身逐渐倾斜而产生的指向内侧的力，也逐渐连续地增大，并应等于火车沿轨道曲线行驶的向心力 $\frac{mv^2}{r}$，这里 $r$ 是曲率半径。所以，为了保证行车安全，火车就不能从直线轨道一下子转入半径为 $R$ 的圆弧轨道，而必须在两者之间插入一段过渡轨道，并且这过渡轨道的曲线形状应该满足一个特殊要求：它的曲率从 $0$ 开始，逐渐均匀地增大到 $\frac{1}{R}$。这里所谓“均匀”，就是指曲率的增大与弧长的增加成正比。

\noindent
\begin{minipage}[t]{0.55\textwidth}
  \vspace{0pt}
  \hspace{2em}曲率 $\varkappa$ 与弧长 $s$ 成正比的曲线，叫做回旋曲线，又叫柯尔努螺线，其形状如图 86.1 所示。图中曲线在坐标原点的曲率为 $0$，向两侧延伸时，曲率随弧长均匀增大，并且分别绕着定点 $M_1$ 和 $M_2$ 旋转。

  \hspace{2em}根据第 85 节关于曲率的公式和第 12 节关于弧长的公式，可以求出回旋曲线的方程。事实上，回旋曲线的定义可写成关系式
  \[
    \varkappa = as,
  \]
\end{minipage}
\hfill
\begin{minipage}[t]{0.42\textwidth}
  \vspace{0pt}
  \centering
  \begin{tikzpicture}[scale=1.35, >=Stealth, line cap=round, line join=round]
    \draw[->] (-1.22, 0) -- (1.22, 0) node[below, font=\footnotesize] {$x$};
    \draw[->] (0, -1.12) -- (0, 1.18) node[left, font=\footnotesize] {$y$};
    \node[below right=-1pt, font=\footnotesize] at (0, 0) {$O$};
    % 精确 Fresnel 积分坐标（1.32*C(s), 1.32*S(s), s in [-3.55, 3.55], 201 点自适应加密, 两端各盘旋约 3.15 圈）
    \draw[line width=0.75pt, curveblue] plot[smooth] coordinates {
      (-0.754,-0.589) (-0.735,-0.568) (-0.712,-0.552) (-0.686,-0.542) (-0.658,-0.539) (-0.629,-0.542)
      (-0.602,-0.552) (-0.578,-0.568) (-0.559,-0.588) (-0.544,-0.614) (-0.536,-0.642) (-0.534,-0.671)
      (-0.539,-0.700) (-0.551,-0.727) (-0.568,-0.751) (-0.591,-0.771) (-0.618,-0.784) (-0.648,-0.791)
      (-0.678,-0.792) (-0.708,-0.785) (-0.736,-0.772) (-0.760,-0.752) (-0.779,-0.727) (-0.793,-0.699)
      (-0.799,-0.668) (-0.798,-0.636) (-0.790,-0.605) (-0.776,-0.576) (-0.755,-0.552) (-0.728,-0.532)
      (-0.698,-0.519) (-0.666,-0.513) (-0.632,-0.514) (-0.600,-0.522) (-0.570,-0.538) (-0.544,-0.560)
      (-0.524,-0.588) (-0.510,-0.619) (-0.503,-0.653) (-0.504,-0.688) (-0.512,-0.723) (-0.528,-0.755)
      (-0.551,-0.782) (-0.579,-0.805) (-0.612,-0.821) (-0.648,-0.829) (-0.685,-0.829) (-0.722,-0.822)
      (-0.756,-0.806) (-0.787,-0.784) (-0.812,-0.755) (-0.831,-0.721) (-0.843,-0.683) (-0.846,-0.644)
      (-0.841,-0.604) (-0.827,-0.566) (-0.806,-0.531) (-0.777,-0.501) (-0.743,-0.478) (-0.704,-0.462)
      (-0.662,-0.454) (-0.619,-0.455) (-0.576,-0.465) (-0.536,-0.483) (-0.500,-0.509) (-0.469,-0.543)
      (-0.446,-0.582) (-0.431,-0.626) (-0.424,-0.673) (-0.427,-0.720) (-0.439,-0.767) (-0.460,-0.812)
      (-0.490,-0.852) (-0.527,-0.886) (-0.571,-0.913) (-0.620,-0.932) (-0.673,-0.941) (-0.727,-0.941)
      (-0.781,-0.931) (-0.834,-0.910) (-0.883,-0.881) (-0.927,-0.842) (-0.965,-0.795) (-0.995,-0.741)
      (-1.016,-0.682) (-1.028,-0.618) (-1.029,-0.552) (-1.019,-0.484) (-0.999,-0.417) (-0.967,-0.351)
      (-0.926,-0.288) (-0.874,-0.230) (-0.812,-0.177) (-0.742,-0.131) (-0.663,-0.091) (-0.575,-0.058)
      (-0.479,-0.033) (-0.375,-0.016) (-0.262,-0.005) (-0.138,-0.001) (0.000,0.000) (0.138,0.001)
      (0.262,0.005) (0.375,0.016) (0.479,0.033) (0.575,0.058) (0.663,0.091) (0.742,0.131)
      (0.812,0.177) (0.874,0.230) (0.926,0.288) (0.967,0.351) (0.999,0.417) (1.019,0.484)
      (1.029,0.552) (1.028,0.618) (1.016,0.682) (0.995,0.741) (0.965,0.795) (0.927,0.842)
      (0.883,0.881) (0.834,0.910) (0.781,0.931) (0.727,0.941) (0.673,0.941) (0.620,0.932)
      (0.571,0.913) (0.527,0.886) (0.490,0.852) (0.460,0.812) (0.439,0.767) (0.427,0.720)
      (0.424,0.673) (0.431,0.626) (0.446,0.582) (0.469,0.543) (0.500,0.509) (0.536,0.483)
      (0.576,0.465) (0.619,0.455) (0.662,0.454) (0.704,0.462) (0.743,0.478) (0.777,0.501)
      (0.806,0.531) (0.827,0.566) (0.841,0.604) (0.846,0.644) (0.843,0.683) (0.831,0.721)
      (0.812,0.755) (0.787,0.784) (0.756,0.806) (0.722,0.822) (0.685,0.829) (0.648,0.829)
      (0.612,0.821) (0.579,0.805) (0.551,0.782) (0.528,0.755) (0.512,0.723) (0.504,0.688)
      (0.503,0.653) (0.510,0.619) (0.524,0.588) (0.544,0.560) (0.570,0.538) (0.600,0.522)
      (0.632,0.514) (0.666,0.513) (0.698,0.519) (0.728,0.532) (0.755,0.552) (0.776,0.576)
      (0.790,0.605) (0.798,0.636) (0.799,0.668) (0.793,0.699) (0.779,0.727) (0.760,0.752)
      (0.736,0.772) (0.708,0.785) (0.678,0.792) (0.648,0.791) (0.618,0.784) (0.591,0.771)
      (0.568,0.751) (0.551,0.727) (0.539,0.700) (0.534,0.671) (0.536,0.642) (0.544,0.614)
      (0.559,0.588) (0.578,0.568) (0.602,0.552) (0.629,0.542) (0.658,0.539) (0.686,0.542)
      (0.712,0.552) (0.735,0.568) (0.754,0.589)
    };
    \node[font=\tiny, inner sep=0.3pt] at (0.660, 0.660) {$M_1$};
    \node[font=\tiny, inner sep=0.3pt] at (-0.660, -0.660) {$M_2$};
  \end{tikzpicture}
  \captionof*{figure}{图\ 86.1}
\end{minipage}

\noindent 其中 $a$ 是常数，$a \neq 0$。根据第 85 节，
\[
  \varkappa = \frac{1}{r} = \frac{d\alpha}{ds},
\]
其中 $\alpha$ 是曲线的切线与 $x$ 轴的夹角。代入上述关系式，得到
\[
  \frac{d\alpha}{ds} = as,
\]
积分得
\[
  \alpha = \frac{a}{2} s^2 + C_1,
\]
其中 $C_1$ 为积分常数。选择坐标系，使得当 $s = 0$ 时 $\alpha = 0$，于是 $C_1 = 0$，有
\[
  \alpha = \frac{as^2}{2}.
\]

另一方面，利用第 12 节的微分三角形可得
\begin{gather*}
  \frac{dx}{ds} = \cos\alpha, \quad \frac{dy}{ds} = \sin\alpha, \\
  \therefore\quad dx = \cos\frac{as^2}{2}\dd s, \quad dy = \sin\frac{as^2}{2}\dd s.
\end{gather*}
积分，并且取 $s = 0$ 时 $x = y = 0$，得到
\[
  x = \int_0^s \cos\frac{as^2}{2}\dd s, \quad y = \int_0^s \sin\frac{as^2}{2}\dd s.
\]
这就是回旋曲线的参数方程。

在铁路或公路转弯处选择过渡曲线时，只用到 $s$ 很小的部分。利用第 84 节的方法，可知当 $s$ 接近于 $0$ 时，有
\[
  \cos\frac{as^2}{2} \approx 1 - \frac{1}{2}\left(\frac{as^2}{2}\right)^2, \quad
  \sin\frac{as^2}{2} \approx \frac{as^2}{2},
\]
因而近似地得到
\begin{align*}
  x &= \int_0^s \left(1 - \frac{a^2}{8} s^4\right)\dd s = s - \frac{a^2}{40} s^5, \\
  y &= \int_0^s \frac{a}{2} s^2\dd s = \frac{a}{6} s^3.
\end{align*}
我国目前的铁路曲线表中，过渡曲线所用的参数方程正是具有形式
\[
  x = t - \frac{a^2}{40} t^5, \quad y = \frac{a}{6} t^3,
\]
其中系数 $\frac{a^2}{40}$ 和 $\frac{a}{6}$ 都已给出具体数值。

如果把 $t^5$ 项也弃掉，消去参数 $t$ 后得到
\[
  y = \frac{a}{6} x^3,
\]
这是一条立方抛物线，有时也用它做过渡曲线。

回旋曲线是物理学家柯尔努在 1874 年研究光的衍射现象时发现的，所以人们也常常把这种螺旋形的曲线叫做柯尔努螺线。
